\section{A Regional SPIN Architecture}
\label{sec:structured-spin}
\label{sec:regional-spin}

We now specify the components of $E=I\circ\Pi\circ O$ from
\eqref{eq:intro-spin-composition}. Split each outer output into packets.
The route sends one packet from every outer group to each region, then
shuffles the packets within each region. The inner processes the regions
in order while carrying a state between them. Packet width, local outer
maps, and inner maps remain parameters here. Sections~\ref{sec:structured-distance}
and~\ref{sec:turbo-spin} fix them for the two distance guarantees.

\subsection{Parameters and Dataflow}

Choose $L$ outer groups, each with $k$ input bits and $b$ output bits.
Split each output into packets of $p\mid b$ bits. The inner processes
$t$ bits per step and holds $s$ state bits. Require $t\mid pL$, so a
region contains whole steps. Put
\begin{equation}
 B:=b/p,\qquad K:=Lk,\qquad N:=Lb,\qquad
 e:=pL/t,\qquad R:=N/t=Be.
 \label{eq:structured-parameters}
\end{equation}
Each of the $B$ regions contains one $p$-bit packet from every group and
occupies $pL$ consecutive output coordinates. Refining
\eqref{eq:intro-spin-composition}, the encoder has the dataflow
\begin{equation}
 E=I\circ\Pi_{\mathrm{reg}}\circ T\circ\Pi_{\mathrm{blk}}\circ O
 :\F_2^{Lk}\longrightarrow\F_2^{Lb}.
 \label{eq:structured-dataflow}
\end{equation}
Here $O$ is a direct sum of local encoders, $\Pi_{\mathrm{blk}}$ shuffles
whole packets within each group, $T$ transposes the group and region
indices, and $\Pi_{\mathrm{reg}}$ shuffles packet slots within each region.
We identify the outer array with $\F_2^N$ in group-major order and the
routed array with $\F_2^N$ in region-major order; $T$ includes this change
of serialization. The three middle maps compose to one permutation of the
$N$ binary coordinates.

The setup consists of injective linear local maps
$G_j:\F_2^k\to\F_2^b$, linear maps $A:\F_2^s\to\F_2^t$ and
$C:\F_2^t\to\F_2^s$, linear state maps $M_i:\F_2^s\to\F_2^s$,
and the route permutations below. An instantiation specifies which maps
are fixed and the joint distribution of the sampled maps. It may reuse
one sampled local map or sample independent maps across groups. Outer,
route, and inner setup are independent in both instantiations. Once
sampled, one setup is fixed for every message.

\subsection{Block-Diagonal Outer}

Write a message as $u=(u_1,\ldots,u_L)$ with
$u_j\in\F_2^k$. The outer encoder is the direct sum
\begin{equation}
  O(u)
  :=
  \bigl(G_1(u_1),\ldots,G_L(u_L)\bigr)
  \in(\F_2^b)^L.
  \label{eq:structured-outer}
\end{equation}
Equivalently, a generator matrix for $O$ places the local generator matrices
on the diagonal. The choice between one shared constituent and independent
local maps affects the word counts in a distance proof. Sections~\ref{sec:structured-distance}
and~\ref{sec:turbo-spin} specify those choices after defining the route.

The first-moment proof needs counts of outer words in the classes relevant
to the route and inner. For one-bit packets, these are row-weight classes;
for four-bit packets, they retain packet support and labels. The two
specializations give their respective counting bounds below.
\subsection{Packet Routing}

Index the outer packets by $(j,r)\in[L]\times[B]$, with packet bit
$\ell\in[p]$. For each group $j$, sample $\sigma_j\gets S_B$ and send
packet $(j,r)$ to region $\sigma_j(r)$ without changing its $p$ bits.
The transpose exchanges the group and region indices. For each region
$r$, independently sample $\pi_r\gets S_L$ and move the packet from
group $j$ to slot $\pi_r(j)$. Thus
\begin{equation}
 \Pi:=\Pi_{\mathrm{reg}}\circ T\circ\Pi_{\mathrm{blk}}.
 \label{eq:structured-interleaver}
\end{equation}
Serialize the result in region-major, then slot-major order, retaining
the bit order within each packet. Region $r$ occupies coordinates
$(r-1)pL+1,\ldots,rpL$.

Every outer group contributes exactly one packet to every region. Given
the set of its nonzero packets, its group shuffle makes their destination
regions a uniform subset of the $B$ regions. Given the packet values
arriving in one region, its regional shuffle gives uniform slot positions.
The route preserves all within-packet structure. At $p=1$, a group of
bit weight $w$ reaches exactly $w$ distinct regions, with one bit per
region. At $p=4$, the same statement concerns nonzero packets, and the
proof must also account for their four-bit values.

\begin{lemma}[Packet-route symmetry]
\label{lem:packet-route-symmetry}
Fix two outer words whose $p$-bit packet multisets agree in every group.
Their routed words have the same distribution. Consequently their
conditional failure probabilities $q_d$ from
\eqref{eq:conditional-inner-failure} agree for every $d$.
\end{lemma}

\begin{proof}
Within each group, one word is a packet permutation of the other.
Composing that permutation with an independent uniform $\sigma_j$ leaves
the group shuffle uniform. The region shuffles and inner setup have the
same independent distributions in both cases.
\end{proof}

\subsection{Recursive Inner and Transpose}

The inner processes the serialized vector in $R$ steps of $t$ bits.
Step inputs $X_i$, outputs $Y_i$, and states $Q_i$ are column vectors.
For fixed maps $A:\F_2^s\to\F_2^t$ and
$C:\F_2^t\to\F_2^s$, set
\begin{equation}
 Q_0:=0,\qquad
 Y_i:=X_i+A Q_i,\qquad
 Q_{i+1}:=M_iQ_i+C X_i
 \quad(0\le i<R).
 \label{eq:structured-inner}
\end{equation}
The state continues across region boundaries; $Q_R$ is discarded.
The recurrence allows any specified distribution for the $M_i$ and any
fixed pair $A,C$. The two main instantiations use $C=A^\top$ and
independent uniform invertible $M_i$. Their choices of $A$ are given
with their distance claims.

The transpose processes input groups $U_i\in\F_2^t$ in reverse order.
Starting from $r_R=0\in\F_2^s$, it emits
\begin{equation}
 V_i=U_i+C^\top r_{i+1},\qquad
 r_i=A^\top U_i+M_i^\top r_{i+1}.
 \label{eq:structured-inner-transpose}
\end{equation}
The feedback computations consume the raw input, not the emitted output.

\begin{lemma}[Inner bijectivity]
\label{lem:structured-inner-bijective}
For every fixed state-map schedule and fixed pair $A,C$, the map
$I:(\F_2^t)^R\to(\F_2^t)^R$ defined in
\eqref{eq:structured-inner} is bijective.
\end{lemma}

\begin{proof}
Fix $Y_0,\ldots,Y_{R-1}$. The initial state is known. At step $i$,
recover $X_i=Y_i+A Q_i$ and then compute $Q_{i+1}$ from
\eqref{eq:structured-inner}. This recovers exactly one input sequence
from every output sequence.
\end{proof}

\subsection{Composition and Rate}

\begin{definition}[Regional SPIN]
\label{def:structured-spin}
For admissible parameters and a setup sampled as specified above,
Regional SPIN is the linear map in \eqref{eq:structured-dataflow}, from
$\F_2^{K}$ to $\F_2^{N}$. A family of parameter choices indexed by $m$
gives maps $E_m:\F_2^{K_m}\to\F_2^{N_m}$.
\end{definition}

Every realized local map is injective. The outer direct sum is therefore
injective. The interleaver is a permutation, and
Lemma~\ref{lem:structured-inner-bijective} makes the inner bijective. Hence
every realized encoder is injective and has rate
\begin{equation}
  \frac{K_m}{N_m}=\frac{k_m}{b_m}.
  \label{eq:structured-spin-rate}
\end{equation}
A distance bound additionally requires outer-class multiplicities and
conditional route--inner failure envelopes as inputs to
Theorem~\ref{thm:class-first-moment}. Lemma~\ref{lem:packet-route-symmetry}
gives an exact class map through the packet multiset of each group;
the proofs may retain these classes or bound unions of them. Scalable SPIN
uses bit-level row weights and an asymptotic envelope. Turbo SPIN retains
packet information for sparse occupations and uses a pointwise outer
comparison for larger ones. The supporting BCH--IMT certificates in
Appendix~\ref{sec:finite-certificates} concern the application backend.
